% vim: set tw=78 sts=2 sw=2 ts=8 aw et:
\documentclass{beamer}

\usepackage[utf8x]{inputenc}		% diacritice
\usepackage[romanian]{babel}
%\usepackage{color}			% highlight
%\usepackage{alltt}			% highlight
%\usepackage{code/highlight}		% highlight
\usepackage{hyperref}			% folosiți \url{http://...}
                                        % sau \href{http://...}{Nume Link}
\usepackage{verbatim}
\usepackage{tabularx}

\mode<presentation>
\usetheme{CDL}

% Titlul nu foloseşte Unicode pentru că e o problemă căreia nu i-am dat de
% cap.
\title[Issue tracking]{Issue tracking}
\subtitle{CDL - Cursul 5}
\institute[ROSEdu]{ROSEdu}
\date{27 martie 2010}
\author{Răzvan Deaconescu \\ \texttt{razvan@rosedu.org}}

\begin{document}

% Slide-urile cu mai multe părți sunt marcate cu textul (cont.)
\setbeamertemplate{frametitle continuation}[from second]
% Arătăm numărul frame-ului
\setbeamertemplate{footline}[frame number]

\frame{\titlepage}

\frame{\tableofcontents}


% NB: Secțiunile nu sunt marcate vizual, ci doar apar în cuprins.
\section{Issue tracking}

% Pentru reamintirea periodică a cuprinsului și unde ne aflăm:
\frame{\tableofcontents[currentsection]}

% Titlul unui frame se specifică fie în acolade, imediat după \begin{frame},
% fie folosind \frametitle

\begin{frame}{Issues}
  \begin{itemize}
    \item probleme, solicitări de rezolvare a unor probleme
    \item tip problemă, autor, asignat, stare, deadline, prioritate,
descriere
    \item tichete -- help desk, call center
    \item issue tracking system -- gestiunea issue-urilor unei
organizații, unui proiect
      \begin{itemize}
        \item aplicație software
        \item interfață web, bază de date
        \item asemănător cu un bug tracking system (bugtracker)
        \item autentificare -- în proiectele open-source submiterea de
bug-uri e deschisă
      \end{itemize}
  \end{itemize}
\end{frame}

\begin{frame}{Issue tracking systems}
  \begin{itemize}
    \item Bugzilla, MantisBT -- single purpose
    \item Trac, Redmine -- multi purpose
    \item SourceForge, Launchpad, Google Code, GitHub -- hosted
  \end{itemize}
\end{frame}


\section{Issue tracking în Redmine}

% Pentru reamintirea periodică a cuprinsului și unde ne aflăm:
\frame{\tableofcontents[currentsection]}

\begin{frame}{New issues}
  \begin{itemize}
    \item tracker -- tipul issue-ului (configurabil)
    \item subject
    \item description
    \item category (de configurat)
    \item due date
    \item priority
    \item assigned to
    \item attached files
  \end{itemize}
\end{frame}

\begin{frame}{View/edit issues}
  \begin{itemize}
    \item list
    \item summary
    \item Gantt chart
    \item calendar
    \item filters
    \item update
    \item delete, move, copy (admin rights)
    \item personalizare mod de afișare
  \end{itemize}
\end{frame}

\section{Concluzii}

% Pentru reamintirea periodică a cuprinsului și unde ne aflăm:
\frame{\tableofcontents[currentsection]}

\begin{frame}{Link-uri utile}
  \begin{itemize}
    \item \url{http://en.wikipedia.org/wiki/Issue\_tracking\_system}
    \item \url{http://www.redmine.org/wiki/redmine/RedmineIssues}
    \item \url{http://www.bugzilla.org/}
    \item
\url{http://en.wikipedia.org/wiki/Comparison\_of\_issue\_tracking\_systems}
    \item \url{http://www.chiark.greenend.org.uk/~sgtatham/bugs.html}
  \end{itemize}
\end{frame}

\begin{frame}{Cuvinte cheie \& întrebări}
  \begin{columns}
    \begin{column}[l]{0.5\textwidth}
      \begin{itemize}
        \item issues
        \item issue tracking
        \item bug tracking
        \item Redmine
        \item create, view, update
      \end{itemize}
    \end{column}
    \begin{column}[c]{0.5\textwidth}
      \begin{figure}
        \includegraphics[scale=0.4]{img/question-mark.jpg}
      \end{figure}
    \end{column}
  \end{columns}
\end{frame}

\begin{frame}{Exerciții}
  \footnotesize \url{http://swarm.cs.pub.ro/~razvan/dokuwiki/rosedu/cdl/issues/readme}
\end{frame}


\end{document}
